\section{The uniform coefficient transseries}
\label{sec:main-proof}

We now prove the main result. The definition of $B$ in
\eqref{eq:scale} incorporates all logarithmic changes of the quadratic
coefficient, including the slowly varying feedback contribution.

\begin{theorem}[Every fixed logarithmic order, uniformly on the lattice]
\label{thm:all-orders}
Fix $\Lambda<\infty$, $0<b_-<b_+<\infty$, and an integer $N\ge0$.
For $L=\log n$, $|\eps L/2|\le\Lambda$, and $b\in[b_-,b_+]$,
let $X_n$, $B$, $q$, and $\xi_m$ be defined by
\eqref{eq:pgf}, \eqref{eq:scale}, and \eqref{eq:q-xi}. Then
\begin{equation}\label{eq:uniform-lattice}
 \sup_{m\in\mathbb Z}\left|
 B\Prob(X_n=m)-\sum_{j=0}^N\mathcal C_j(\xi_m;q,\eps)
 \right|\le C_N L^{-N-1}
\end{equation}
for all sufficiently large $n$, uniformly in the stated parameters.
Consequently the feedback coefficients satisfy
\begin{equation}\label{eq:all-orders-coefficient}
 \boxed{\displaystyle
 nB R_b^n a_n
 =\normal(\xi)+
  \sum_{j=1}^N\frac1{j!}\sum_{k=1}^j
    \stir{j}{k}\left(-\frac q2\right)^k\eps^{j-k}
           \HH_{2k}^{[j]}(\xi)
       +O(L^{-N-1}).}
\end{equation}
The finite divided-difference expression
\eqref{eq:finite-difference-approx} has the same error bound.
\end{theorem}

The theorem resolves the joint coefficient target Q1 in
\cite{repo-boundary}. It is an absolute estimate. For bounded $\xi$
it gives relative Gaussian asymptotics because $\normal$ has a
positive minimum on each compact real interval. Relative conclusions
for unbounded $\xi$ will be proved separately.

\subsection{Uniform control of the whole Fourier contour}

\begin{lemma}[Uniform Fourier domination]
\label{lem:fourier-tail}
Let
\begin{equation}\label{eq:chi}
 \chi_n(y)=\E\exp\left(iy\frac{X_n-nA/b}{B}\right).
\end{equation}
There are positive constants $c_1,c_2$, uniform in the moving-order
window and in $b$, such that
\begin{align}
 |\chi_n(y)|&\le e^{-c_1y^2}
       &&(|y|\le Bn^{-1/4}),\label{eq:fourier-small}\\
 |\chi_n(y)|&\le e^{-c_2\sqrt n}
       &&(Bn^{-1/4}\le |y|\le\pi B).
       \label{eq:fourier-large}
\end{align}
\end{lemma}

\begin{proof}
By positivity of the intensities,
\[
 -\log|\chi_n(B\theta)|
  =n\sum_{k\ge1}c_k(1-\cos k\theta).
\]
For $|\theta|\le n^{-1/4}$ retain $1\le k\le\lfloor n^{1/4}\rfloor$.
Then $|k\theta|\le1$, so $1-\cos k\theta\ge c k^2\theta^2$.
The parameters keep $u$ bounded above. Moreover
$k^{-\eps}\ge e^{-\Lambda/2}$ on this range, since
$|\eps|\log k\le\Lambda/2$. Therefore
\[
 -\log|\chi_n(B\theta)|
 \ge c n\theta^2\sum_{k\le n^{1/4}}k^{-1-\eps}
 \ge c' nL\theta^2.
\]
Since $B^2\asymp nL$, this proves \eqref{eq:fourier-small}.
For $n^{-1/4}\le|\theta|\le\pi$, retain only $k=1$.
Here $c_1=1/u$ in the intensity notation and
$1-\cos\theta\ge2\theta^2/\pi^2$, giving
$-\log|\chi_n(B\theta)|\ge c n\theta^2\ge c\sqrt n$.
The centering phase has modulus one, proving the second assertion.
\end{proof}

This argument is insensitive to the cancellation of the Gamma and
zeta poles. It uses the exact nonnegative coefficients and supplies
the uniformity that a fixed-order local-limit theorem alone would
not provide.

\subsection{The exact central reduction}

For $z\in\C$, define the entire divided exponential
\begin{equation}\label{eq:h-epsilon}
 h_\eps(z)=\frac{e^{\eps z}-1}{\eps},\qquad h_0(z)=z.
\end{equation}
Take $t=-iy/B$ in \cref{lem:local}. The scale equation removes the
constant quadratic term and gives
\begin{equation}\label{eq:core-characteristic}
 \chi_n(y)=\exp\left\{
     -\frac{y^2}{2}+\frac q2y^2h_\eps(\Log(-iy))+R_n(y)
                       \right\},
 \quad |R_n(y)|\le CnB^{-5/2}|y|^{5/2},
\end{equation}
whenever $0<|y|<rB$. The limit at zero is understood continuously.
The sign in this formula follows from
$(-iy)^2=-y^2$ and
\[
 K_\eps(\log B-\Log(-iy))-K_\eps(\log B)
   =-E_2(\eps)B^{-\eps}h_\eps(\Log(-iy)).
\]
Since $nB^{-5/2}=O(n^{-1/4}L^{-5/4})$, the local nonlinear
remainder is smaller than every fixed inverse power of $L$.

\subsection{A finite expansion with an integrable remainder}

\begin{proof}[Proof of \cref{thm:all-orders}]
Choose
$T_L=\sqrt{D_N\log L}$ with a sufficiently large constant $D_N$.
By \cref{lem:fourier-tail}, the part of the scaled Fourier integral
with $T_L<|y|\le\pi B$ has absolute size
$O(L^{-N-1})$: use the Gaussian bound up to $Bn^{-1/4}$ and
the bound $2\pi B e^{-c_2\sqrt n}$ beyond it.

On $|y|\le T_L$, put $z=\Log(-iy)$ and
$w=(q/2)y^2h_\eps(z)$. The inequality
\begin{equation}\label{eq:h-bound}
 |h_\eps(z)|\le |z|e^{|\eps z|}
\end{equation}
follows from $h_\eps(z)=z\int_0^1e^{v\eps z}\,dv$.
It proves $\sup_{|y|\le T_L}|w|=o(1)$. Near $y=0$ one uses
$y^2|\log|y||\,|y|^{-|\eps|}\to0$; no uniform bound on
$|\eps\log|y||$ is assumed. For $1\le|y|\le T_L$ the bound is
$O(L^{-1}\log L\log\log L)$.
Equation \eqref{eq:core-characteristic} and the estimate on $R_n$
therefore allow us to replace $\chi_n$ by
$e^{-y^2/2}e^w$ in the central integral with error
$O(nB^{-5/2})$.

Expand the exponential only to the finite order $N$:
\[
 e^w=\sum_{k=0}^N\frac{w^k}{k!}+O_N(|w|^{N+1}).
\]
The error is integrable against $e^{-y^2/2}$ with bound
$C_N q^{N+1}$. Indeed \eqref{eq:h-bound} bounds the integrand near
zero by a constant times
$|y|^{2N+2-(N+1)|\eps|}(1+|\log|y||)^{N+1}$;
at infinity a Gaussian dominates the polynomial and logarithmic
factors. For the fixed $N$, take $n$ large enough that
$(N+1)|\eps|<1/2$. All these integrals then have a common majorant.

For each $k\le N$, use
\begin{equation}\label{eq:cube-integral}
 h_\eps(z)^k=z^k\int_{[0,1]^k}
      e^{\eps z(v_1+\cdots+v_k)}\,dv_1\cdots dv_k.
\end{equation}
Taylor expansion in $\eps$ through degree $N-k$ gives
\[
 h_\eps(z)^k
 =k!\sum_{j=k}^N\stir{j}{k}\frac{\eps^{j-k}z^j}{j!}
    +\mathcal R_{k,N}(z),
\]
where
\[
 |\mathcal R_{k,N}(z)|\le
 C_N|\eps|^{N+1-k}|z|^{N+1}e^{k|\eps z|}.
\]
After multiplication by $q^k|y|^{2k}e^{-y^2/2}$, its integral is
$O(q^k|\eps|^{N+1-k})=O(L^{-N-1})$, by the same domination.
Thus the total-degree polynomial in $q,\eps$ has an $L^1(dy)$
error of the asserted order. Its finite terms may be extended from
$|y|\le T_L$ to the full real axis: their Gaussian tails are
$O(L^{-N-1})$ after enlarging $D_N$ if needed.

The resulting integrals are
\[
 \frac1{2\pi}\int_\R e^{-i\xi_m y-y^2/2}
       (-iy)^{2k}(\Log(-iy))^j\,dy
       =\HH_{2k}^{[j]}(\xi_m).
\]
Here the factor $y^{2k}$ in $w^k$ contributes $(-1)^k$, producing
$(-q/2)^k$. Lattice Fourier inversion says
\[
 B\Prob(X_n=m)
 =\frac1{2\pi}\int_{-\pi B}^{\pi B}
                 e^{-i\xi_m y}\chi_n(y)\,dy.
\]
Every error bound above is an $L^1$ bound before insertion of the
phase. Since $|e^{-i\xi_m y}|=1$, it is uniform in all integers $m$.
This proves \eqref{eq:uniform-lattice}. The exact coefficient identity
gives \eqref{eq:all-orders-coefficient}.

If one retains $h_\eps(z)^k$ instead of Taylor-expanding in
$\eps$, apply finite Fourier inversion to
\[
 (-iy)^{2k}h_\eps(\Log(-iy))^k.
\]
This gives precisely \eqref{eq:finite-difference-approx}. The same
$q^{N+1}$ estimate proves its assertion.
\end{proof}

\begin{remark}[The full formal exponential is not the Fourier law]
The proof integrates a finite expansion in the central region and
uses the exact characteristic function on its complement. It does not
replace the full Fourier law by
$\exp\{-y^2/2+(q/2)y^2h_\eps(\Log(-iy))\}$ on the whole real axis.
For fixed positive $q,\eps$ the latter expression can grow at large
$|y|$. Nor does the theorem assert convergence of the infinite sum
over $j$. The finite integrals and all fixed-order remainders are the
claimed objects.
\end{remark}

\subsection{The elementary normalization requested in Q1}

Put $\xi_0=n(1-A/b)/B_0$, using \eqref{eq:leading-scale}.
On compact $\xi_0$ sets, \cref{thm:scale,thm:all-orders} give
\begin{equation}\label{eq:elementary-first}
\begin{split}
 nB_0R_b^n a_n={}&\normal(\xi_0)
   +\frac{C_\lambda\HH_2(\xi_0)
               -e^{-\lambda}\HH_2^{[1]}(\xi_0)}{L\ph(\lambda)}\\
 &+O\left(\frac{(\log L)^2}{L^2}\right).
\end{split}
\end{equation}
To verify the scale-change term, write
$B/B_0=1+\delta$; then
$(B_0/B)\normal(\xi_0B_0/B)
=\normal(\xi_0)-\delta(\normal+\xi_0\normal')(\xi_0)+O(\delta^2)$,
and $-(\normal+\xi\normal')=\HH_2$.

In particular, if $\xi_0$ remains bounded, the leading statement is
\begin{equation}\label{eq:Q1-solved}
 n\sqrt{\frac{n\log n}{2b}\ph(\lambda)}\,R_b^n a_n
 =\normal\left(
 \frac{n(1-\zeta(2+\eps)/b)}{
               \sqrt{n\log n\,\ph(\lambda)/(2b)}}\right)+o(1).
\end{equation}
At the critical choice $b=A=\zeta(2+\eps)$,
\begin{equation}\label{eq:critical-joint-leading}
 a_n\sim R_b^{-n}
   \sqrt{\frac{A}{\pi\ph(\lambda)}}\,
                   n^{-3/2}(\log n)^{-1/2},
\end{equation}
uniformly over the bounded moving-order window. This is a joint
transition law; letting $\lambda$ tend to infinity is outside its
stated uniform parameter set.
